% !TEX root = ../main.tex
\chapter{Mathlib の CategoryTheory へ進む}
\label{chap:mathlib-category-theory}

\begin{chapterabstract}
本章では、本書で使ってきた自作のミニ定義から、Mathlib の \texttt{CategoryTheory} ライブラリへ進むための最初の地図を作る。目的は、Mathlib の圏論 API を網羅することではない。\texttt{Category}、\texttt{Functor}、\texttt{NatTrans}、\texttt{Iso} という入口を読み、本書で学んだ「対象・射・合成・同型・関手・自然変換」が、Mathlib ではどの名前と記法で現れるかを対応づけることである。

ここまで本書では、概念の理解を優先して小さな構造体を自作してきた。本格的な形式化に進むときは、その自作定義を捨てるのではなく、Mathlib の語彙へ翻訳する。本章は、その翻訳のための足場である。
\end{chapterabstract}

\begin{learninggoals}
\item 自作ミニ定義と Mathlib の \texttt{CategoryTheory} 定義の役割の違いを説明できる。
\item \texttt{Category}、\texttt{Functor}、\texttt{NatTrans}、\texttt{Iso} の基本フィールドを読める。
\item \(X \longrightarrow Y\)、\(\mathbb{1}_X\)、\texttt{f ≫ g}、\(C \Rightarrow D\)、\(X \cong Y\) の記法を本書の言葉へ戻せる。
\item いきなり Mathlib 圏論から始める難しさを、型クラス・universe・暗黙引数の観点から理解できる。
\item 自作定義を Mathlib の語彙へ移すための小さな手順を持てる。
\end{learninggoals}

\section{この章を読む理由}

ここまでの章では、圏論の概念を実務に近い形で導入してきた。圏は合成できる処理の世界であり、同型は情報を失わない相互変換であり、関手は構造を保った一括変換であり、自然変換は一貫した adapter であった。Lean コードでも、これらを小さな \texttt{structure} と関数で表した。

しかし、Lean 上で既存の数学的結果を使いたい場合や、他の形式化プロジェクトと接続したい場合、自作定義だけでは足りない。Mathlib にはすでに圏論の大きなライブラリがある。そこには圏、関手、自然変換、同型、極限、随伴、モナド、具体圏、加法圏、アーベル圏など、多くの定義と定理が蓄積されている。

本章の目的は、その巨大なライブラリに正面から飛び込むことではない。まず「本書で見た概念が Mathlib のどこにあるか」を確認する。そして、Mathlib の記法を見たときに、実務寄りの意味へ戻して読めるようにする。

\begin{intuitionbox}
自作ミニ定義は、学習用の模型である。模型は小さく、見通しがよく、どこに何が入っているかを自分で確認できる。一方、Mathlib の定義は、実プロジェクトで再利用するための標準部品である。標準部品は強力だが、型クラス、記法、暗黙引数、universe が絡むため、最初は見通しが悪い。本章では、模型から標準部品へ乗り換えるための対応表を作る。
\end{intuitionbox}

\FigurePlaceholder{fig-ch34-overview.png}{0.90}{自作ミニ定義から Mathlib CategoryTheory へ進む地図}{fig:ch34-overview}

\section{本書の自作ミニ定義と Mathlib の違い}

本書の前半では、説明しやすさを優先して、圏や関手を小さな構造体として作った。たとえば、対象の型、射の型、恒等射、合成、法則をすべて一つの構造体に入れるような定義である。この形は、圏論の考え方を学ぶには向いている。

Mathlib の定義は、それよりも Lean の型クラス機構に強く合わせて設計されている。特に重要なのは、\textbf{対象の集まりをフィールドとして持つ構造体ではなく、対象の型 \texttt{C} に対して \texttt{Category C} という型クラスを載せる}という点である。

\begin{softwarebox}
ソフトウェア的に言えば、自作ミニ定義は「設定オブジェクトに全部を詰める」設計に近い。Mathlib の \texttt{Category C} は、「型 \texttt{C} に、圏として振る舞うためのインターフェース実装を与える」設計に近い。つまり、対象は \texttt{C} の値であり、射・恒等射・合成・法則は型クラスから取り出される。
\end{softwarebox}

\begin{longtable}{p{0.25\linewidth}p{0.30\linewidth}p{0.35\linewidth}}
\toprule
本書のミニ定義での言葉 & Mathlib でよく見る形 & 読み替え \\
\midrule
対象 & \texttt{X : C} & 圏 \texttt{C} の中のオブジェクト。型 \texttt{C} の値として表される。 \\
射 & \(X \longrightarrow Y\) & \texttt{X} から \texttt{Y} への morphism。ソフトウェアでは変換・処理・対応として読む。 \\
恒等射 & \(\mathbb{1}_X\) & 対象 \texttt{X} 上の「何もしない射」。 \\
合成 & \texttt{f ≫ g} & 先に \texttt{f}、次に \texttt{g} を適用する矢印向きの合成。 \\
圏 & \texttt{[Category C]} & 型 \texttt{C} に圏構造がある、という型クラス引数。 \\
関手 & \(F : C \Rightarrow D\) & 圏 \texttt{C} から圏 \texttt{D} への構造保存写像。 \\
対象への作用 & \texttt{F.obj X} & 関手が対象をどう移すか。 \\
射への作用 & \texttt{F.map f} & 関手が射をどう移すか。 \\
自然変換 & \texttt{NatTrans F G} & 関手 \texttt{F} から \texttt{G} への一貫した変換。 \\
自然変換の成分 & \texttt{α.app X} & 各対象 \texttt{X} での adapter。 \\
同型 & \(X \cong Y\) & 互いに戻せる射を束ねたもの。 \\
\bottomrule
\end{longtable}

\begin{warningbox}
Lean の標準ライブラリにも \texttt{Functor} や \texttt{Monad} という名前がある。これはプログラミング用の型クラスであり、\texttt{List.map} や \texttt{Option.bind} のような計算上の抽象化を扱う。一方、Mathlib の \texttt{CategoryTheory.Functor} は一般の圏の間の関手である。名前は似ているが、扱う抽象度が違う。本章では常に \texttt{CategoryTheory} 名前空間の圏論用語を扱う。
\end{warningbox}

\section{\texttt{Category}：対象の型に圏構造を載せる}

Mathlib で圏を読むとき、まず見るべき形は次である。

\begin{listing}[htbp]
\begin{minted}{lean4}
import Mathlib.CategoryTheory.Category.Basic

universe v u

namespace Chapter34

open CategoryTheory
open scoped CategoryTheory

section CategoryFields

variable {C : Type u} [Category.{v} C]
variable {W X Y Z : C}

#check (X ⟶ Y)
#check (𝟙 X)
#check fun (f : X ⟶ Y) (g : Y ⟶ Z) => f ≫ g

end CategoryFields
end Chapter34
\end{minted}
\caption{圏構造を仮定して射・恒等射・合成を見る}
\label{lst:ch34_category_basic}
\end{listing}

ここで \texttt{C} は対象の型である。\texttt{[Category.\{v\} C]} は、「\texttt{C} の値を対象とする圏構造を型クラス推論で使う」という意味で読む。射の型は \(X \longrightarrow Y\) と書く。恒等射は \(\mathbb{1}_X\)、合成は \texttt{f ≫ g} と書く。

\begin{definitionbox}
Mathlib の \texttt{Category C} は、対象型 \texttt{C} に対して、射 \(X \longrightarrow Y\)、恒等射 \(\mathbb{1}_X\)、合成 \texttt{f ≫ g}、および恒等律・結合律を与える型クラスである。本書の第5章で作った「対象・射・合成・恒等射・法則を持つもの」は、Mathlib ではこの \texttt{Category C} として表される。
\end{definitionbox}

圏の法則も、Mathlib ではすでに \texttt{simp} で扱いやすい形になっている。

\begin{listing}[htbp]
\begin{minted}{lean4}
section CategoryLaws

variable {C : Type u} [Category.{v} C]
variable {W X Y Z : C}

theorem id_then (f : X ⟶ Y) : (𝟙 X) ≫ f = f := by
  simpa

theorem then_id (f : X ⟶ Y) : f ≫ (𝟙 Y) = f := by
  simpa

theorem comp_assoc_example
    (f : W ⟶ X) (g : X ⟶ Y) (h : Y ⟶ Z) :
    (f ≫ g) ≫ h = f ≫ (g ≫ h) := by
  simpa using (Category.assoc f g h)

end CategoryLaws
\end{minted}
\caption{恒等律と結合律を Mathlib の記法で使う}
\label{lst:ch34_category_laws}
\end{listing}

\begin{leanbox}
\texttt{simpa} は、\texttt{simp} による整理と、現在のゴールへの適合を同時に行う道具である。ここでは、圏の恒等律が単純化規則として登録されているため、恒等射を前後に合成した式が自動的に整理される。結合律については、\texttt{Category.assoc} を明示して使っている。
\end{leanbox}

\section{\texttt{Functor}：対象と射を同時に運ぶ}

本書では、\texttt{List.map} や \texttt{Option.map} を使って、「構造を保ったまま中身を変える」操作として関手を導入した。Mathlib の一般圏論では、関手は圏から圏への構造保存写像である。記法は \(C \Rightarrow D\) である。

\begin{listing}[htbp]
\begin{minted}{lean4}
import Mathlib.CategoryTheory.Functor.Basic

universe v1 v2 u1 u2

namespace Chapter34

open CategoryTheory
open scoped CategoryTheory

section FunctorFields

variable {C : Type u1} [Category.{v1} C]
variable {D : Type u2} [Category.{v2} D]
variable (F : C ⥤ D)

#check F.obj
#check F.map
#check fun {X Y : C} (f : X ⟶ Y) => F.map f

end FunctorFields
end Chapter34
\end{minted}
\caption{Functor の obj と map を読む}
\label{lst:ch34_functor_basic}
\end{listing}

\texttt{F.obj X} は、関手 \texttt{F} が対象 \texttt{X} をどの対象へ移すかを表す。\texttt{F.map f} は、射 \texttt{f} をどの射へ移すかを表す。ここまでは、本書で作ったミニ関手とほぼ同じ感覚で読める。

違いは、Mathlib の関手が一般の圏を対象にしている点である。\texttt{List.map} のような「型の中身を変える操作」だけでなく、群の圏、位相空間の圏、加法圏、関手圏など、さまざまな圏の間の構造保存写像を同じ型で扱う。

\begin{listing}[htbp]
\begin{minted}{lean4}
section FunctorLaws

variable {C : Type u1} [Category.{v1} C]
variable {D : Type u2} [Category.{v2} D]
variable (F : C ⥤ D)

theorem functor_preserves_id (X : C) :
    F.map (𝟙 X) = 𝟙 (F.obj X) := by
  simpa using F.map_id X

theorem functor_preserves_comp {X Y Z : C}
    (f : X ⟶ Y) (g : Y ⟶ Z) :
    F.map (f ≫ g) = F.map f ≫ F.map g := by
  simpa using F.map_comp f g

end FunctorLaws
\end{minted}
\caption{関手が恒等射と合成を保存することを使う}
\label{lst:ch34_functor_laws}
\end{listing}

\begin{softwarebox}
実務で \texttt{F.map} を読むときは、「ある変換を、別の文脈へ壊さずに持ち上げる」と考えるとよい。第24章のバッチマイグレーションでは、1件の \texttt{migrate} を \texttt{List.map migrate} でリスト全体へ持ち上げた。Mathlib の関手では、その考え方を任意の圏の間へ一般化している。
\end{softwarebox}

\section{\texttt{NatTrans}：各対象の成分と自然性}

自然変換は、二つの関手の間の変換である。本書では、API response wrapper の adapter として読んだ。Mathlib では、自然変換 \texttt{α : NatTrans F G} は、各対象 \texttt{X} に対する成分 \texttt{α.app X} と、任意の射に対する自然性の等式を持つ。

\begin{listing}[htbp]
\begin{minted}{lean4}
import Mathlib.CategoryTheory.NatTrans

universe v1 v2 u1 u2

namespace Chapter34

open CategoryTheory
open scoped CategoryTheory

section NatTransFields

variable {C : Type u1} [Category.{v1} C]
variable {D : Type u2} [Category.{v2} D]
variable {F G : C ⥤ D}
variable (α : NatTrans F G)

#check fun (X : C) => α.app X
#check fun {X Y : C} (f : X ⟶ Y) => α.naturality f

end NatTransFields
end Chapter34
\end{minted}
\caption{NatTrans の app と naturality を読む}
\label{lst:ch34_nattrans_basic}
\end{listing}

自然性は、次のような等式として現れる。

\begin{listing}[htbp]
\begin{minted}{lean4}
section NatTransNaturality

variable {C : Type u1} [Category.{v1} C]
variable {D : Type u2} [Category.{v2} D]
variable {F G : C ⥤ D}
variable (α : NatTrans F G)

theorem naturality_example {X Y : C} (f : X ⟶ Y) :
    F.map f ≫ α.app Y = α.app X ≫ G.map f := by
  simpa using α.naturality f

end NatTransNaturality
\end{minted}
\caption{自然性の四角形を等式として使う}
\label{lst:ch34_nattrans_naturality}
\end{listing}

この式は、第9章と第25章で見た可換図式そのものである。左辺は「先に \texttt{F} 側で処理してから adapter を通す」経路であり、右辺は「先に adapter を通してから \texttt{G} 側で処理する」経路である。両者が等しいとき、adapter は業務ロジックと干渉しない。

\begin{intuitionbox}
\texttt{α.app X} は、型 \texttt{X} ごとに用意された adapter である。しかし、それだけでは「一貫した adapter」とは言えない。自然性 \texttt{α.naturality f} があることで、どの処理 \texttt{f} と組み合わせても順序を入れ替えられる、という保証が加わる。
\end{intuitionbox}

\section{\texttt{Iso}：round-trip をライブラリの語彙で読む}

第6章では、同型を「情報を失わない相互変換」として導入した。\texttt{to} と \texttt{from} があり、両方向の round-trip law を証明する、という形である。

Mathlib では、圏 \texttt{C} の対象 \texttt{X} と \texttt{Y} の同型は \(X \cong Y\) と書く。中身には、前向きの射 \texttt{hom}、後ろ向きの射 \texttt{inv}、そして両方向の合成が恒等射になる証明が入っている。

\begin{listing}[htbp]
\begin{minted}{lean4}
import Mathlib.CategoryTheory.Iso

universe v u

namespace Chapter34

open CategoryTheory
open scoped CategoryTheory

section IsoFields

variable {C : Type u} [Category.{v} C]
variable {X Y : C}
variable (e : X ≅ Y)

#check e.hom
#check e.inv
#check e.hom_inv_id
#check e.inv_hom_id

end IsoFields
end Chapter34
\end{minted}
\caption{Iso を round-trip として読む}
\label{lst:ch34_iso_basic}
\end{listing}

\begin{listing}[htbp]
\begin{minted}{lean4}
section IsoRoundTrip

variable {C : Type u} [Category.{v} C]
variable {X Y : C}
variable (e : X ≅ Y)

theorem iso_roundtrip_left : e.hom ≫ e.inv = 𝟙 X := by
  simpa using e.hom_inv_id

theorem iso_roundtrip_right : e.inv ≫ e.hom = 𝟙 Y := by
  simpa using e.inv_hom_id

end IsoRoundTrip
\end{minted}
\caption{同型の二つの round-trip law}
\label{lst:ch34_iso_roundtrip}
\end{listing}

\begin{softwarebox}
スキーマ移行に戻して読むなら、\texttt{e.hom} が \texttt{migrate}、\texttt{e.inv} が \texttt{rollback} に対応する。\texttt{e.hom\_inv\_id} は「移行して戻すと元に戻る」、\texttt{e.inv\_hom\_id} は「戻して移行すると元に戻る」という主張である。ただし Mathlib の \texttt{Iso} は一般の圏の射についての同型であり、通常の関数同型だけに限られない。
\end{softwarebox}

\section{いきなり Mathlib 圏論から始めない理由}

Mathlib の \texttt{CategoryTheory} は強力である。しかし、圏論と Lean の両方を学び始めた段階で、いきなりここから入るとつまずきやすい。理由は、概念そのもの以外の学習負荷が大きいからである。

\subsection{型クラス推論が前提になる}

Mathlib では、\texttt{[Category C]} のような型クラス引数が頻繁に出てくる。これは、「\texttt{C} が圏であることを Lean が文脈から探して使う」という仕組みである。便利な一方で、エラーが出たときに「圏構造が見つからない」のか、「射の向きが違う」のか、「universe が合っていない」のかを見分ける必要がある。

\subsection{universe が見える}

小さな自作定義では、対象も射も同じ \texttt{Type} レベルに置いても大きな問題にならない。しかし Mathlib では、型の圏や関手圏のように大きさの違う対象を扱うため、universe level が明示的に現れる。\texttt{Category.\{v\} C} の \texttt{v} は、射の型が住む universe を指定している。

\subsection{記法が多い}

Mathlib の圏論では、\(X \longrightarrow Y\)、\(\mathbb{1}_X\)、\texttt{f ≫ g}、\(C \Rightarrow D\)、\(X \cong Y\) など、読みやすさのための記法が多く使われる。慣れると自然だが、最初はどれが型で、どれが項で、どれが定理なのかが見えにくい。

\begin{warningbox}
Mathlib の記法を見て混乱したときは、記法を本書の言葉へ戻すとよい。\(X \longrightarrow Y\) は「\texttt{X} から \texttt{Y} への射」、\texttt{f ≫ g} は「\texttt{f} の後に \texttt{g}」、\(C \Rightarrow D\) は「\texttt{C} から \texttt{D} への関手」である。記法を覚えることより、何を保存しているかを読むことが重要である。
\end{warningbox}

\FigurePlaceholder{fig-ch34-mini-to-mathlib.png}{0.92}{自作ミニ定義を Mathlib の語彙へ移す対応関係}{fig:ch34-mini-to-mathlib}

\section{発展課題：自作定義を Mathlib の語彙へ移す}

ここでは、実際に移行するときの考え方を整理する。ゴールは、過去の章で作ったミニ定義を一気に置き換えることではない。まず、どの概念が Mathlib のどの構成要素に対応するかを一つずつ確認する。

\subsection{第一段階：名前の対応を作る}

最初にやることは、証明を書き直すことではなく、名前の対応表を作ることである。

\begin{longtable}{p{0.30\linewidth}p{0.30\linewidth}p{0.30\linewidth}}
\toprule
移行前に探すもの & Mathlib で探すもの & 確認すること \\
\midrule
\texttt{MiniCategory.Obj} & \texttt{C : Type u} & 対象は型の値として表す。 \\
\texttt{MiniCategory.X ⟶ Y} & \(X \longrightarrow Y\) & 射の向きが一致しているか。 \\
\texttt{id} & \(\mathbb{1}_X\) & 恒等射の対象が合っているか。 \\
\texttt{comp f g} & \texttt{f ≫ g} & 合成順序が逆になっていないか。 \\
\texttt{MiniFunctor.obj} & \texttt{F.obj} & 対象への作用が一致しているか。 \\
\texttt{MiniFunctor.map} & \texttt{F.map} & 射への作用が一致しているか。 \\
\texttt{MiniNatTrans.app} & \texttt{α.app} & 各対象での成分が一致しているか。 \\
\texttt{MiniIso.to/from} & \texttt{e.hom/e.inv} & round-trip の向きが一致しているか。 \\
\bottomrule
\end{longtable}

\subsection{第二段階：小さな定理だけを移す}

次に、過去の章で使った小さな定理を一つだけ選び、Mathlib の記法で書き直す。おすすめは、次の順である。

\begin{enumerate}
\item 恒等射を合成しても変わらないこと。
\item 関手が恒等射を保存すること。
\item 関手が合成を保存すること。
\item 自然変換の自然性。
\item 同型の round-trip law。
\end{enumerate}

これらは本章のコード例にすでに出てきた。つまり、過去の自作定義で証明した内容は、Mathlib では既存のフィールドや simp ルールとして利用できる場合が多い。

\begin{leanbox}
Mathlib に移すときは、「自分で証明する」よりも「どの既存定理が対応するかを探す」作業が増える。たとえば、関手の合成保存は \texttt{F.map\_comp f g}、自然性は \texttt{α.naturality f}、同型の round-trip は \texttt{e.hom\_inv\_id} と \texttt{e.inv\_hom\_id} としてすでに入っている。
\end{leanbox}

\subsection{第三段階：実務モデルに戻す}

Mathlib の \texttt{CategoryTheory} で抽象定理を読めるようになったら、もう一度、実務のモデルへ戻す。たとえば API adapter の自然性を、一般の \texttt{NatTrans} として完全に表す必要があるかを判断する。多くの実務モデルでは、第25章のような小さな \texttt{map} と adapter の等式で十分である。

本格的な Mathlib 圏論を使うべき場面は、既存の圏論的定理を再利用したいときである。単に業務ロジックの round-trip を証明したいだけなら、自作の小さなモデルのほうが読みやすいことも多い。

\section{実務への接続：いつ Mathlib を使うべきか}

Mathlib の \texttt{CategoryTheory} は、すべての実務証明に必要なわけではない。むしろ、最初の導入では小さな自作モデルのほうがチームに説明しやすい。では、いつ Mathlib へ進むべきか。

\begin{softwarebox}
Mathlib を使う価値が大きいのは、証明したい内容がすでに標準的な数学構造に乗っているときである。たとえば、一般の関手・自然変換・同型・極限・随伴に関する定理を使いたい場合、Mathlib の語彙に乗せることで既存ライブラリの定理を再利用できる。一方、単一システム内の小さなデータ変換や業務仕様だけを扱うなら、自作モデルで十分な場合がある。
\end{softwarebox}

判断基準は次のように置くとよい。

\begin{longtable}{p{0.42\linewidth}p{0.48\linewidth}}
\toprule
状況 & 推奨する方針 \\
\midrule
チーム内で Lean の導入を始める & 自作の小さなモデルで仕様の核を表す。 \\
API adapter や migration の局所性を示す & まずは本書の小さな定義で証明する。 \\
圏論の標準定理を使いたい & Mathlib の \texttt{CategoryTheory} に移す。 \\
他の Mathlib 形式化と接続したい & Mathlib の標準定義を使う。 \\
研究寄りの形式化を進める & 最初から Mathlib の語彙で設計する。 \\
\bottomrule
\end{longtable}

\section{本章のまとめ}

Mathlib の \texttt{Category C} は、対象型 \texttt{C} に圏構造を載せる型クラスとして読む。

\(X \longrightarrow Y\) は射、\(\mathbb{1}_X\) は恒等射、\texttt{f ≫ g} は矢印向きの合成である。

\(F : C \Rightarrow D\) は関手であり、\texttt{F.obj} と \texttt{F.map} によって対象と射を移す。

\texttt{NatTrans F G} は自然変換であり、\texttt{α.app X} と \texttt{α.naturality f} が中心になる。

\(X \cong Y\) は同型であり、\texttt{hom}、\texttt{inv}、二つの round-trip law を持つ。

Mathlib は強力だが、型クラス、universe、記法の負荷があるため、本書のミニ定義で概念を掴んでから進むのが安全である。

実務では、Mathlib を使うこと自体が目的ではない。既存の標準定義と定理を再利用したいときに、必要な範囲で移行する。

\section{参考リンク}

本章の Mathlib API 名は、次の公式ドキュメントを入口として確認できる。

\begin{itemize}
\item \texttt{Mathlib.CategoryTheory.Category.Basic}: \url{https://leanprover-community.github.io/mathlib4_docs/Mathlib/CategoryTheory/Category/Basic.html}
\item \texttt{Mathlib.CategoryTheory.Functor.Basic}: \url{https://leanprover-community.github.io/mathlib4_docs/Mathlib/CategoryTheory/Functor/Basic.html}
\item \texttt{Mathlib.CategoryTheory.NatTrans}: \url{https://leanprover-community.github.io/mathlib4_docs/Mathlib/CategoryTheory/NatTrans.html}
\item \texttt{Mathlib.CategoryTheory.Iso}: \url{https://leanprover-community.github.io/mathlib4_docs/Mathlib/CategoryTheory/Iso.html}
\item Lean のプログラミング用 \texttt{Functor} / \texttt{Monad} との違い: \url{https://lean-lang.org/doc/reference/latest/Functors___-Monads-and--do--Notation/}
\end{itemize}
